\begin{initiativkarte}{Neues Fliegen e.V.}

  % Bild (optional)
  \IfFileExists{bilder/neues_fliegen.jpg}{%
    \begin{center}
      \includegraphics[width=\linewidth,height=5cm,keepaspectratio]{bilder/neues_fliegen.jpg}
    \end{center}
    \vspace{0.4cm}
  }{}

  % Kurzbeschreibung
  \textit{\textcolor{hawgray}{Neues Fliegen e.V. beschäftigt sich mit der Zukunft der Luftfahrt und entwickelt innovative Flugmodelle und Wettbewerbsformate im studentischen Umfeld.}}

  \vspace{0.5cm}
  \hawrule

  % Beschreibung
  \textbf{\textcolor{hawblue}{Was macht Neues Fliegen?}}\\[0.3cm]

  Der studentische Verein \textbf{Neues Fliegen e.V.} beschäftigt sich mit der Zukunft des Fliegens und überträgt aktuelle Fragestellungen der Luftfahrt auf den Modellbau.

  Ziel ist es, Studierenden eine praxisnahe Möglichkeit zu geben, Studieninhalte durch Konstruktion, Entwicklung und Organisation im realen Kontext anzuwenden.

  \vspace{0.3cm}

  Dabei sind Studierende aller Fachrichtungen willkommen – nicht nur aus der Luftfahrttechnik. Die Aufgaben im Verein sind vielfältig und reichen von Konstruktion und Fertigung über IT-Systeme bis hin zu Organisation und Sponsoring.

  \vspace{0.5cm}

  \textbf{New Flying Competition}

  Ein zentrales Projekt ist die \textbf{New Flying Competition}, ein internationaler studentischer Wettbewerb, der 2027 in die 6. Ausgabe startet.

  Teams aus verschiedenen Ländern treten mit selbst entwickelten Lösungen gegen ein Referenzflugzeug von Neues Fliegen an, um eine definierte Aufgabe möglichst effizient zu lösen.

  Die Aufgaben werden in Zusammenarbeit mit Industriepartnern entwickelt und greifen reale Herausforderungen der Luftfahrt auf.

  Bei der letzten Ausgabe nahmen Teams aus acht Ländern teil, darunter Australien, China und Mexiko.

  \vspace{0.5cm}

  Ziel ist es, technische, organisatorische und persönliche Fähigkeiten zu verbinden und Studierenden einen realitätsnahen Einblick in die Luftfahrtentwicklung zu geben.

  \vspace{0.5cm}

  % Tags
  \tag{Luftfahrt}
  \tag{Engineering}
  \tag{Modellbau}
  \tag{Wettbewerb}
  \tag{Innovation}

  \vspace{0.6cm}
  \hawrule

  % Infozeilen
  \infozeile{\faUsers}{Zielgruppe: Alle Studierenden (technik- und nicht-technikorientiert)}
  \infozeile{\faMapMarker*}{Teamraum: BT11-019}
  \infozeile{\faGlobe}{Website: \href{https://neuesfliegen.de/}{neuesfliegen.de}}

\end{initiativkarte}